% ===== 原文 pp.6–8 =====
% 边界说明：按 TIMELINE #1（修订），本块自 p.6 顶部起。开头一段是 §2.1 的续文（正交 Higgs 场的特征多项式，无编号条目），随后是 §2.2 的小节标题。
% [原文 p.6]

考察正交 Higgs 场 $\Phi$ 的结构时可以看到：对一般的矩阵 $A\in\fso(2n+1,\C)$（对 $A\in\fso(2n,\C)$ 亦然），其互不相同的特征值都以 $\pm\lambda_i$ 的形式成对出现。因此，$A$ 的特征多项式——它在后面几章中将变得重要——必形如
\[
\det(x-A)=x(x^{2n}+a_1x^{2n-2}+\cdots+a_{n-1}x^{2}+a_n).
\]

% [Tutor Note: p.6 开头「补出的省略论证」：$\fso(2n+1,\C)$ 中的矩阵是奇数维反对称矩阵，故必有特征值 $0$（$\det A=0$）；而反对称矩阵的非零特征值关于原点对称，即以 $\pm\lambda_i$ 成对出现，于是 $\det(x-A)=x\prod_{i=1}^{n}(x^{2}-\lambda_i^{2})$，展开即得原文所写的 $x(x^{2n}+a_1x^{2n-2}+\cdots+a_n)$。注意：原文把 $A\in\fso(2n+1,\C)$ 与 $A\in\fso(2n,\C)$ 并列，但含 $x$ 因子的这个形式只适用于奇数维；偶数维时 $\det(x-A)$ 不含 $x$ 因子、只含偶次幂。]

\subsection{实 Higgs 丛}

实形式（real form）的 Higgs 丛最早由 N. Hitchin 在 [60] 中研究；关于 $\Sl(2,\R)$ 的结果在 [62] 中得到推广，其中 Hitchin 研究了通过李代数上的对合（involution）而出现的 $G=\Sl(n,\R)$ 情形——这一视角我们将在后面几章中再回来讨论。利用 Higgs 丛，他数出了连通分支的数目；并且在分裂实形式（split real form）的情形，他找出了一个同胚于 $\R^{\dim G(2g-2)}$ 的分支，该分支自然地包含一个 Teichmüller 空间的副本。

沿用 [61]：对 $\GC$ 的非紧实实形式 $G$（实结构记为 $\tau$）而言，为得到相应的 $G$-Higgs 丛，解 Hitchin 方程 (2.7)--(2.8) 的联络 $A$ 需要满足
\[
\nabla=\nabla_A+\Phi-\rho(\Phi),
\]
才能使其和乐（holonomy）落在 $G$ 中；这里 $\rho$ 是 $\GC$ 的紧实结构（compact real structure）。

由于联络 $A$ 的和乐落在 $\GC$ 的紧实形式中，故有 $\rho(\nabla_A)=\nabla_A$。于是，要求 $\nabla=\tau(\nabla)$ 等价于要求 $\nabla_A=\tau(\nabla_A)$ 与 $\Phi-\rho(\Phi)=\tau(\Phi-\rho(\Phi))$。若用对合 $\sigma:=\rho\tau$ 来表达，这两个等式即为
\begin{gather*}
\sigma(\nabla_A)=\nabla_A\\
\sigma(\rho(\Phi)-\Phi)=\Phi-\rho(\Phi).
\end{gather*}

由上可知：若 $\nabla_A$ 在 $\sigma$ 下不变、而 $\Phi$ 反不变，则 $\nabla$ 的和乐落在实形式 $G$ 中。因此，设 $G$ 是复半单李群 $\GC$ 的实形式，我们可以如下构造 $G$-Higgs 丛。设 $H$ 是 $G$ 的极大紧子群；我们已经看到，$\fg$ 的 Cartan 分解为 $\fg=\fh\oplus\fm$，其中 $\fh$ 是 $H$ 的李代数，$\fm$ 是其正交补。这诱导出复李代数的分解 $\fg_{\C}=\fh_{\C}\oplus\fm_{\C}$。通过诱导出的各向同性表示（isotropy representation）$\mathrm{Ad}|_{H_{\C}}:H^{\C}\to\GL(\fm^{\C})$，便得到 $G$-Higgs 丛的具体描述（更多细节，例如可参见 [87]）：

% [Tutor Note: p.6 首现「各向同性表示」（isotropy representation）：$\mathrm{Ad}|_{H_{\C}}:H^{\C}\to\GL(\fm^{\C})$ 是 $H^{\C}$ 在 $\fm^{\C}$ 上的伴随作用。顺带指出，原文同一句内 $\mathrm{Ad}|_{H_{\C}}$ 用下标、而后面的 $H^{\C}$、$\GL(\fm^{\C})$ 用上标，上下标混用（p.7 亦作 $H^{\C}$）；本译本照原文排印，读者可统一读作 $H^{\C}$。]

\begin{defn}{2.13} 主 $G$-Higgs 丛（principal $G$-Higgs bundle）是一对 $(P,\Phi)$，其中 $P$ 是 $\Sigma$ 上的全纯主 $H^{\C}$-丛，而 $\Phi$ 是丛 $P\times_{\mathrm{Ad}}\fm^{\C}\otimes K$ 的全纯截面。\end{defn}

% [Tutor Note: p.6 定义 2.13 中，相伴丛 $P\times_{\mathrm{Ad}}\fm^{\C}$ 的纤维是 $\fm^{\C}$，结构群为 $H^{\C}$；因此 $\Phi$ 在每点取值于 $\fm^{\C}\otimes K$，这正是把 $G$-Higgs 丛写成「主丛 + 截面」数据的关键。]

\begin{exmp}{2.14} 考虑 $\Sl(p+q,\C)$ 的非紧实形式 $\SU(p,q)$，其李代数为
\[
\fsu(p,q)=\left\{\begin{pmatrix}Z_1&Z_2\\ \overline{Z}_2^{\,t}&Z_3\end{pmatrix}\right\},
\]
其中 $Z_1,Z_3$ 分别是 $p$ 阶与 $q$ 阶的斜 Hermitian 矩阵，$\Tr Z_1+\Tr Z_3=0$，而 $Z_2$ 任意。

% [Tutor Note: p.6 例 2.14 的矩阵元 $\overline{Z}_2^{\,t}$ 是 $Z_2$ 的共轭转置；文字层会把它抽成 $Z_2^t$（丢掉上划线），易误读为普通转置。可代 $p=q=1$ 校验维数：$Z_1,Z_3$ 的斜 Hermitian 条件给出 $p^2+q^2$ 个实参数，$Z_2$ 任意给出 $2pq$ 个，再减去 $\Tr Z_1+\Tr Z_3=0$ 一个约束，得 $(p+q)^2-1$，与 $\dim\SU(p+q)$ 相符。]

% [原文 p.7]

其复化具有 Cartan 分解
\[
\fsu(p,q)^{\C}=\fsl(p+q,\C)=\fh^{\C}\oplus\fm^{\C},
\]
其中 $\fm^{\C}$ 对应于 $\fsl(p+q,\C)$ 的非对角元素。$\SU(p,q)$ 的中心是 $\U(1)$，其极大紧子群为 $H=S(\U(p)\times\U(q))$，而复化后的李群为
\[
H^{\C}=\{(X,Y)\in\GL(p,\C)\times\GL(q,\C): \det Y=(\det X)^{-1}\}.
\]
因此，$\Sigma$ 上的 $\SU(p,q)$-Higgs 丛是一对 $(E,\Phi)$，其中 $E=V_p\oplus V_q$，$V_p,V_q$ 是 $\Sigma$ 上秩分别为 $p$、$q$ 的向量丛，满足 $\Lambda^{p}V_p\cong\Lambda^{q}V_q^{*}$，而 Higgs 场 $\Phi$ 由
\[
\Phi=\begin{pmatrix}0&\beta\\ \gamma&0\end{pmatrix}
\]
给出，其中 $\beta: V_q\to V_p\otimes K$，$\gamma: V_p\to V_q\otimes K$。\end{exmp}

% [Tutor Note: p.7 的条件 $\Lambda^{p}V_p\cong\Lambda^{q}V_q^{*}$ 等价于 $\det V_p\otimes\det V_q\cong\mathcal{O}_{\Sigma}$，正是把结构群约化到 $S(\U(p)\times\U(q))$ 的条件（与紧邻的 $H^{\C}$ 定义 $\det Y=(\det X)^{-1}$ 相对照）。]

关于实 Higgs 丛的构造及其几何性质的研究，更多细节读者或许愿意参阅讲义与学位论文 [7, 51, 82, 89, 93]，以及下面这份——当然绝非详尽——我们过去觉得有用的论文清单（及其中的参考文献）：[8,11,14,28,34,35,46,48,49,62,64,73,86,87,91,92,101,104]。特别地，关于实形式的 Hitchin--Kobayashi 型对应，读者应当参阅 [47]。

\subsection{抛物 Higgs 丛}

Higgs 丛的概念可以在若干方向上推广；过去数十年间，其中一个引起数学家和物理学家注意的方向是抛物 Higgs 丛（parabolic Higgs bundle），它源于向量丛向抛物丛（parabolic bundle）的推广 [74]。

在亏格为 $g$、带有标记点 $p_1,\dots,p_n$ 且满足 $2g-2+n\ge 0$ 的紧黎曼面 $\Sigma$ 上，记 $D$ 为有效除子 $D=p_1+\dots+p_n$。于是按 [74]，$\Sigma$ 上的抛物向量丛 $E$ 是指 $\Sigma$ 上的一个代数秩 $r$ 向量丛 $E$，再带上一个抛物结构（parabolic structure）——即各标记点处 $E$ 的纤维上的一个（未必完整的）旗：
\[
0=E_{p_i,0}\subset E_{p_i,1}\subset E_{p_i,2}\subset\cdots\subset E_{p_i,r_i}=E|_{p_i},
\]
并且对每个标记点 $p_i$ 配上一组抛物权（parabolic weight）$\alpha_i=(\alpha_{i,0},\dots,\alpha_{i,r_i})$；这些权是满足
\[
1=\alpha_{i,0}>\alpha_{i,1}>\cdots>\alpha_{i,r_i}\ge 0
\]
的实数。

% [Tutor Note: p.7 的抛物权约定 $1=\alpha_{i,0}>\alpha_{i,1}>\cdots>\alpha_{i,r_i}\ge0$ 比通常「$0\le\alpha<1$」的约定多出一项 $\alpha_{i,0}=1$；而 p.8 的抛物次数求和是 $\sum_{i=1}^{r_j}$（不含 $i=0$），故这一项只是规范化约定、不参与求和，两种约定给出的抛物次数相同。]

抛物全纯映射（parabolic holomorphic map）是指抛物丛之间的映射 $f: E\to E'$，它满足：只要 $\alpha_i(p_j)>\alpha'_k(p_j)$，就有 $f(E_{p_j,i})\subset E'_{p_j,k+1}$（对一切 $p_j\in D$），其中 $\alpha'_k(p_j)$ 是 $E'$ 上的权。抛物丛 $E$ 的抛物次数（parabolic degree）则定义为
% [原文 p.8]
\[
\operatorname{par}\deg(E)=\deg(E)+\sum_{j=1}^{n}\sum_{i=1}^{r_j}m_{j,i}\alpha_{j,i},
\]
其中 $m_{j,i}=\dim(E_{p_j,i}/E_{p_j,i+1})$；而抛物丛的抛物斜率（parabolic slope）定义为 $\operatorname{par}\mu(E)=\operatorname{par}\deg(E)/\rk(E)$。称一个抛物丛是（半）稳定的，如果对 $E$ 的每个抛物子丛 $F$，抛物斜率都满足 $\operatorname{par}\mu(F)\le\operatorname{par}\mu(E)$（相应地，$\operatorname{par}\mu(F)<\operatorname{par}\mu(E)$）。借助 $E$ 的抛物稳定性，可以构造半稳定抛物向量丛的模空间 [74]，通常记作 $\mathcal{N}_{\alpha}$；它是正规射影簇，当权为一般（generic）时非奇异。

% [Tutor Note: p.8 的 $m_{j,i}=\dim(E_{p_j,i}/E_{p_j,i+1})$ 与 p.7 的递增旗 $0=E_{p_j,0}\subset E_{p_j,1}\subset\cdots\subset E_{p_j,r_j}=E|_{p_j}$ 不完全自洽：商 $E_{p_j,i}/E_{p_j,i+1}$ 要求 $E_{p_j,i+1}\subset E_{p_j,i}$（即降旗），而在递增旗下该商无定义。按递增旗，第 $i$ 个逐次商应为 $E_{p_j,i}/E_{p_j,i-1}$（其中 $E_{p_j,0}=0$），这样 $\sum_i m_{j,i}=\rk E$ 才成立。本译本照原文排；读者按「第 $i$ 层的重数 = 该层与其前一层的维数差」理解即可。]

人们可以推广抛物向量丛的概念，考虑抛物 Higgs 丛（parabolic Higgs bundle，简称 PHG），这最早由 C. Simpson 在 [33] 中给出。

\begin{defn}{2.15} （经典）抛物 Higgs 丛（classical parabolic Higgs bundle）是一对 $(E,\Phi)$，其中
\begin{itemize}[leftmargin=1.4em,itemsep=1pt,topsep=1pt]
\item $E$ 是如上定义的抛物向量丛，
\item Higgs 场 $\Phi\in H^{0}(\Sigma,\End(E)\otimes K_{\Sigma}(D))$ 满足 $\Phi|_{p_i}(E_{i,j})\subset E_{i,j}\otimes K_{\Sigma}(D)|_{p_i}$。
\end{itemize}
\end{defn}

稳定性的概念可以推广到抛物 Higgs 丛；由此得到次数为 $d$、秩为 $r$ 的 $\alpha$-半稳定抛物 Higgs 丛的模空间 $\mathcal{P}_{\alpha}$，它最早由 Yokogawa 在 [105] 中构造（另见 [27]）。该空间的维数为 $(2g-2)r^{2}+2+nr(r-1)$，是正规拟射影簇；[70] 用规范场论把它构造为超凯勒商（hyperkähler quotient），而秩 $r=2$ 的抛物 Higgs 丛模空间的拓扑性质，则由 Hitchin 在 [105] 中用 Morse 理论方法研究。

% [Tutor Note: p.8 原文 "This space hasdimension …" 中 hasdimension 连写、缺空格，应为 "has dimension"（500 dpi 页图与文字层均复核为连写）。维数公式可用 $n=0$ 校验：退化为 Hitchin 丛模空间的 $(2g-2)r^{2}+2$。]

% [Tutor Note: p.8 的维数 $(2g-2)r^{2}+2+nr(r-1)$ 恰是半稳定抛物向量丛模空间维数 $(g-1)r^{2}+1+\tfrac{n}{2}r(r-1)$ 的两倍——抛物 Higgs 丛模空间是它的（抛物）余切丛。以 $g=2,r=2,n=1$ 代入自检：后者为 $4+1+1=6$，乘 2 得 12，与 $(2g-2)r^{2}+2+nr(r-1)=8+2+2=12$ 一致。]

\begin{exmp}{2.16} $\U(p,q)$-抛物 Higgs 丛（$U(p,q)$-parabolic Higgs bundle）是 $\Sigma$ 上的抛物 Higgs 丛 $(E,\Phi)$，其中 $E=V\oplus W$，$V$ 与 $W$ 分别是秩 $p$ 与秩 $q$ 的抛物向量丛，且
\[
\Phi=\begin{pmatrix}0&\beta\\ \gamma&0\end{pmatrix} : (V\oplus W)\to(V\oplus W)\otimes K(D),
\]
其中非零分量 $\beta: W\to V\otimes K(D)$ 与 $\gamma: V\to W\otimes K(D)$ 是抛物态射，而 $D$ 是 $\Sigma$ 上的有效除子（读者应与例 2.14 比较）。\end{exmp}

抛物 Hitchin 系统可以看成允许规范场与 Higgs 场在黎曼面的标记点处有单极点；而更高阶的极点则引向野 Higgs 丛（wild Higgs bundle），这是我们接下来要讨论的对象。

\subsection{野 Higgs 丛}

在整份讲义中，我们只限于研究射影直线 $\PP^{1}$ 上的野 Higgs 丛 $(E,\Phi)$。对 $\PP^{1}$ 中固定的有效除子 $D=p_1+\dots+p_n$，我们把这样的对看成由向量丛 $E$ 与一个取值于 $\mathcal{O}_{\PP^{1}}(D)$ 的自同态 $\Phi: E\to E\otimes\mathcal{O}_{\PP^{1}}(D)$（称为 Higgs 场）组成。除子 $D$ 控制 $\Phi$ 允许出现极点的位置，因此它有时被称为极点除子（polar divisor）。
